\chapter{Transformaciones: tono y tiempo}\label{cap:transformaciones}

Para tocar un sonido en distintas teclas hay que cambiar su altura. Hay dos maneras, y
corresponden a los dos modos del sampler de Fourierizer. \emph{Tape} cambia la velocidad: el tono y
la duración cambian juntos, como en una cinta. \emph{Spectral} cambia el tono y conserva la
duración. Este capítulo explica las dos, y por qué la segunda es mucho más difícil.

\section{Tape: remuestrear}

Transponer $s$ semitonos al estilo cinta es leer la señal a otra velocidad. Con
$r=2^{s/12}$, la salida es
\[
  y[j]=x(jr),
\]
donde $x(\cdot)$ es la señal ``continua'' que interpolan las muestras. Doce semitonos arriba,
$r=2$: se lee el doble de rápido, todo suena una octava arriba y dura la mitad.

El teorema de muestreo (\cref{cap:fourier}) dice cómo interpolar sin inventar nada: con la función
$\mathrm{sinc}(u)=\sin(\pi u)/(\pi u)$,
\[
  x(\tau)=\sum_{n}x[n]\,\mathrm{sinc}(\tau-n).
\]
Si $r>1$, leer más rápido sube todas las frecuencias por un factor $r$, y lo que estaba por encima
de $\Fs/(2r)$ rebasaría Nyquist y se \emph{doblaría} (aliasing). Por eso, al acortar, la sinc se
ensancha por $r$ para filtrar antes a la nueva Nyquist:
\[
  y[j]=f_c\sum_n x[n]\,\mathrm{sinc}\bigl(f_c(jr-n)\bigr),\qquad f_c=\min(1,1/r).
\]

La suma infinita se trunca con una ventana de Kaiser (\cref{cap:ventanas}) de $24$ cruces por cero
por lado, con $\beta=8$, y el núcleo se lee de una tabla de $512$ puntos por cruce con interpolación
lineal~\cite{smith2011}. Es el mismo remuestreador que convierte los archivos de $44.1$ a $48$ kHz al
importarlos, y el que usa el \emph{Pitch} de cada capa.

\begin{estudio}
\textbf{Tape tiene carácter.} Las notas agudas son más cortas y más ``apretadas'', y las graves más
lentas y oscuras. Es el sonido de los samplers clásicos y suele ser lo mejor para baterías y
golpes: el ataque se conserva intacto, solo que más rápido o más lento.
\end{estudio}

\section{Spectral: cambiar el tono sin cambiar la duración}

Fourierizer lo hace en dos pasos, los dos sin aliasing:
\begin{enumerate}
  \item \textbf{Estirar} la señal $r$ veces en el tiempo \emph{sin cambiar su tono}, con un
        \emph{phase vocoder}.
  \item \textbf{Remuestrear} por $r$ como en Tape. Leer la señal estirada $r$ veces más rápido
        regresa la duración original y multiplica todas las frecuencias por $r$.
\end{enumerate}
Lo difícil es el primer paso.

\subsection{El phase vocoder}

Para estirar, se analiza la señal con fragmentos separados por un salto $H_a$ y se resintetiza con
fragmentos separados por $H_s=rH_a$. Si solo se cambiaran las posiciones, cada parcial volvería a
empezar su fase en cada fragmento y las sumas se cancelarían: el resultado sonaría como un
temblor. La fase de cada coeficiente tiene que \emph{avanzar} entre fragmentos de síntesis lo que
avanzaría la parcial en $H_s$ muestras:
\[
  \phi^{\text{s}}_{j}(k)=\phi^{\text{s}}_{j-1}(k)+H_s\,\omega_k,
\]
donde $\omega_k$ es la \emph{frecuencia instantánea} de lo que suena en el bin $k$, en radianes
por muestra~\cite{flanagan1966}. El bin solo dice que la frecuencia está \emph{cerca} de
$\theta_k=2\pi k/N$. La precisa se mide con la fase.

\begin{proposicion}[Frecuencia instantánea]\label{prop:frecuencia-instantanea}
Sea $x[n]=A\cos(\theta_0n+\phi)$ y sean $\chi_a(k)$ y $\chi_{a-\delta}(k)$ los coeficientes de dos
fragmentos analizados en las posiciones $a$ y $a-\delta$. Si $|\theta_0-\theta_k|<\pi/\delta$ y la
contribución de la frecuencia negativa es despreciable, entonces
\[
  \theta_0=\theta_k+\frac{1}{\delta}\,\mathrm{princarg}\Bigl(\angle\chi_a(k)-\angle\chi_{a-\delta}(k)-\theta_k\delta\Bigr),
\]
donde $\mathrm{princarg}$ lleva un ángulo a $[-\pi,\pi)$.
\end{proposicion}

\begin{proof}
Por el \cref{cor:parcial}, $\chi_a(k)\approx\tfrac A2\,\e^{\iu(\theta_0 a'+\phi)}\,W(\theta_k-\theta_0)$, donde
$a'$ es el inicio del fragmento. El factor $W$ es el mismo para los dos fragmentos (misma ventana,
mismo bin), así que la diferencia de fases es $\theta_0\delta$ módulo $2\pi$. Es decir,
$\theta_0\delta=\Delta\phi+2\pi j$ para algún entero $j$. Escribimos
$\theta_0\delta=\theta_k\delta+(\theta_0-\theta_k)\delta$: como $|(\theta_0-\theta_k)\delta|<\pi$,
el término $(\theta_0-\theta_k)\delta$ es justamente el representante en $[-\pi,\pi)$ de
$\Delta\phi-\theta_k\delta$. Despejando se obtiene la fórmula.
\end{proof}

En los vocoders clásicos, $\delta$ es el propio salto de análisis. Fourierizer mide la frecuencia
con un segundo fragmento a $\delta=N/8$ muestras, \emph{independiente} de $H_a$. Así la
estimación es válida en $|\theta_0-\theta_k|<\pi/\delta$, es decir, a menos de $N/(2\delta)=4$ bins
del centro, lo que cubre el lóbulo principal de Hann. El salto de análisis puede ser cualquiera,
incluso mayor que la ventana, como ocurre al transponer hacia abajo.

\subsection{Phasiness y phase locking}

Si cada bin avanza su fase por su cuenta, los bins del lóbulo principal de una misma parcial
pierden la relación de fase que tenían en el análisis. La parcial se ``despega'' de sí misma y el
resultado suena difuso y metálico: es la \emph{phasiness}. Laroche y Dolson propusieron
\emph{amarrar} las fases~\cite{laroche1999}:
\begin{enumerate}
  \item Se buscan los picos de la magnitud, es decir, los bins mayores que sus dos vecinos de cada
        lado.
  \item Solo los picos propagan su fase con la frecuencia instantánea.
  \item Cada bin gira \emph{el mismo ángulo} que el pico más cercano. Si el pico pasa de
        $\angle\chi$ a $\phi^{\text{s}}$, todos los bins de su región se multiplican por
        $\e^{\iu(\phi^{\text{s}}-\angle\chi_{\text{pico}})}$.
\end{enumerate}
Así, el lóbulo principal de cada parcial se mueve en bloque y conserva su forma.

En estéreo, Fourierizer calcula esas rotaciones una sola vez, sobre la señal media $L+R$, y aplica
la misma a los dos canales. Como cada bin de cada canal gira lo mismo, las diferencias de fase entre
canales (la imagen estéreo) se conservan.

\subsection{Transientes}

El phase vocoder estira bien lo que es estable. Un ataque, en cambio, cae dentro de varios
fragmentos de síntesis que se superponen, y se reparte en el tiempo: aparece un \emph{pre-eco} antes
del golpe y el ataque pierde filo. El fenómeno es del orden de la ventana. En un tren de clicks
transpuesto $+7$ semitonos, el pre-eco medido fue de $-21\dB$ con $N=4096$, $-28\dB$ con $2048$ y
$-39\dB$ con $1024$. Por eso la ventana del modo Spectral es una elección del usuario: \emph{Tonal}
($4096$), \emph{Mixed} ($2048$) o \emph{Punchy} ($1024$). Una alternativa que se evaluó, reiniciar
las fases en los ataques, apenas mejoró la medición, así que se descartó.

\section{Del algoritmo al instrumento}

Un phase vocoder de alta calidad es demasiado costoso para correr nota por nota en tiempo real.
Fourierizer aprovecha que el sonido a tocar (el resultado) ya está fijo. Calcula una copia por
semitono de transposición, empezando por las notas que tocas y luego la octava alrededor del
\emph{Root}, en dos hilos de fondo y con un presupuesto de memoria. Cada copia se produce por
pedazos y se puede tocar mientras se termina: una nota nueva espera solo su primer pedazo, unos
10 ms. El cálculo por pedazos da exactamente los mismos números que el cálculo completo, y los
tres motores coinciden con la referencia en Python (\cref{cap:instrumento}).

\begin{estudio}
\textbf{Tape o Spectral.}
\begin{itemize}
  \item \emph{Tape} para baterías, golpes y todo lo que deba conservar el ataque intacto, o cuando
        quieres el carácter de sampler clásico.
  \item \emph{Spectral} para pads, texturas, voces y acordes: un pad de 20 s dura 20 s en cualquier
        tecla. Con \emph{Tonal} los sostenidos quedan más limpios; con \emph{Punchy}, los ataques.
  \item Lejos del Root (más de una octava), los dos modos se notan: Tape por la velocidad y
        Spectral por un timbre algo más artificial. Elegir bien el Root ayuda más que cualquier
        ajuste.
\end{itemize}
\end{estudio}
